% Autorska priprava k Realisticky.cz 020204; ulohy jsou autorske a obmenene.
% Zamerne nenahrazuje vyklad presnymi kopiemi cizich textu/prikladu.
\documentclass[11pt,a4paper]{article}
\input{../../sablony/pripravy/priprava.sty}
\usepackage{tikz}
\nastavPripravu{1. ročník}{Rovnice a nerovnice}{75}{Lineární rovnice II: zlomky}
\begin{document}
\blok{Návaznost a motivace}
Na posledních hodinách jsme vyřešili rovnice s celými koeficienty. Teď se přidají zlomky, ale pravidlo zůstává stejné: násobíme celou rovnici vhodným nenulovým společným násobkem jmenovatelů.
\blok{Princip a odvození}
V rovnici s konstantními nenulovými jmenovateli je násobení jejich nejmenším společným násobkem ekvivalentní. Pokud se ve jmenovateli objeví neznámá, nejdříve určíme podmínky a vyloučíme nulový jmenovatel.
\blok{Řešené příklady: od jednoduchého ke složitějšímu}
\begin{prikladReseniBox}[title={Příklad 1}]
Vyřešte $\frac{x-1}{3}+\frac{x+2}{2}=4$.
\tcblower
\begin{align*}
2(x-1)+3(x+2)&=24 \\
 2x-2+3x+6&=24 \\
 5x&=20 \\
 x&=4
\end{align*}

Zkouška: $3/3+6/2=4$.
\end{prikladReseniBox}
\begin{prikladReseniBox}[title={Příklad 2}]
Vyřešte $\frac{3x-2}{4}-\frac{x+1}{3}=1$.
\tcblower
\begin{align*}
3(3x-2)-4(x+1)&=12 \\
 9x-6-4x-4&=12 \\
 5x&=22 \\
 x&=\frac{22}{5}
\end{align*}
\end{prikladReseniBox}
\poznamka{Kontroluj násobení celé rovnice, nikoli jednotlivých vybraných členů. Zkoušku prováděj v původním zlomkovém zápisu.}
\blok{Samostatné procvičování}
\begin{zadaniReseniBox}[title={Úloha 1}]
$\frac{x+1}{2}=\frac{2x-1}{3}$.
\tcblower
\begin{align*}
3(x+1)&=2(2x-1) \\
 3x+3&=4x-2 \\
 x&=5
\end{align*}
\end{zadaniReseniBox}
\begin{zadaniReseniBox}[title={Úloha 2}]
$2(x-3)-3(x+1)=5-x$.
\tcblower
\begin{align*}
-x-9&=5-x \\
 -9&=5 \\
 K&=\varnothing
\end{align*}
\end{zadaniReseniBox}
\begin{zadaniReseniBox}[title={Úloha 3}]
$\frac{x-2}{5}+\frac{x+3}{10}=1$.
\tcblower
\begin{align*}
2(x-2)+(x+3)&=10 \\
 3x-1&=10 \\
 x&=\frac{11}{3}
\end{align*}
\end{zadaniReseniBox}
\blok{Volitelná rezerva}
\begin{zadaniReseniBox}[title={Úloha 1}]
Vysvětlete, proč nelze při řešení $\frac{x+2}{x-1}=2$ násobit bez podmínky.
\tcblower
\begin{align*}
x-1&\ne0\Rightarrow x\ne1 \\
 x+2&=2(x-1) \\
 x&=4
\end{align*}

Kořen $4$ splňuje podmínku. Bez uvedení $x\ne1$ není řešení úplné.
\end{zadaniReseniBox}
\blok{Domácí úkol}
\begin{zadaniReseniBox}[title={Úloha 1}]
$\frac{2x+1}{3}=\frac{x+4}{2}$.
\tcblower
\begin{align*}
2(2x+1)&=3(x+4) \\
 4x+2&=3x+12 \\
 x&=10
\end{align*}
\end{zadaniReseniBox}
\begin{zadaniReseniBox}[title={Úloha 2}]
$1-\frac{x-2}{4}=\frac{x+2}{2}$.
\tcblower
\begin{align*}
4-(x-2)&=2(x+2) \\
 6-x&=2x+4 \\
 x&=\frac23
\end{align*}
\end{zadaniReseniBox}
\begin{zadaniReseniBox}[title={Úloha 3}]
$4x+2=4x+2$.
\tcblower
\begin{align*}
4x+2&=4x+2 \\
 0&=0 \\
 K&=\mathbb R
\end{align*}
\end{zadaniReseniBox}
\poznamka{Orientační časové rozvržení: 5 + 12 + 12 + 10 + 6 = 45 min. Přednost má pochopení jádra; rezerva není součástí časového součtu. Domácí úlohy jsou určeny na samostatnou práci, řešení zůstávají učiteli.}
\end{document}
